% Options for packages loaded elsewhere
\PassOptionsToPackage{unicode}{hyperref}
\PassOptionsToPackage{hyphens}{url}
\documentclass[
  spanish,
  12pt,
  letterpaper,
]{article}
\usepackage{xcolor}
\usepackage[margin=2.4cm]{geometry}
\usepackage{amsmath,amssymb}
\setcounter{secnumdepth}{-\maxdimen} % remove section numbering
\usepackage{iftex}
\ifPDFTeX
  \usepackage[T1]{fontenc}
  \usepackage[utf8]{inputenc}
  \usepackage{textcomp} % provide euro and other symbols
\else % if luatex or xetex
  \usepackage{unicode-math} % this also loads fontspec
  \defaultfontfeatures{Scale=MatchLowercase}
  \defaultfontfeatures[\rmfamily]{Ligatures=TeX,Scale=1}
\fi
\usepackage{lmodern}
\ifPDFTeX\else
  % xetex/luatex font selection
  \setmainfont[]{Times New Roman}
\fi
% Use upquote if available, for straight quotes in verbatim environments
\IfFileExists{upquote.sty}{\usepackage{upquote}}{}
\IfFileExists{microtype.sty}{% use microtype if available
  \usepackage[]{microtype}
  \UseMicrotypeSet[protrusion]{basicmath} % disable protrusion for tt fonts
}{}
\usepackage{setspace}
\makeatletter
\@ifundefined{KOMAClassName}{% if non-KOMA class
  \IfFileExists{parskip.sty}{%
    \usepackage{parskip}
  }{% else
    \setlength{\parindent}{0pt}
    \setlength{\parskip}{6pt plus 2pt minus 1pt}}
}{% if KOMA class
  \KOMAoptions{parskip=half}}
\makeatother
\ifLuaTeX
\usepackage[bidi=basic,shorthands=off,]{babel}
\else
\usepackage[bidi=default,shorthands=off,]{babel}
\fi
\ifPDFTeX
\else
\babelfont{rm}[]{Times New Roman}
\fi
\ifLuaTeX
  \usepackage{selnolig} % disable illegal ligatures
\fi
\setlength{\emergencystretch}{3em} % prevent overfull lines
\providecommand{\tightlist}{%
  \setlength{\itemsep}{0pt}\setlength{\parskip}{0pt}}

\usepackage{titlesec}
\titleformat{\section}{\centering\bfseries}{\thesection}{1em}{}
\titleformat{\subsection}{\raggedright\bfseries}{\thesubsection}{1em}{}
\titleformat{\subsubsection}{\raggedright\bfseries\itshape}{\thesubsubsection}{1em}{}
\titleformat{\paragraph}[runin]{\raggedright\bfseries}{\theparagraph}{1em}{}[.]
\titleformat{\subparagraph}[runin]{\raggedright\bfseries\itshape}{\thesubparagraph}{1em}{}
\titlespacing*{\paragraph}{0cm}{*1}{1em}
\titlespacing*{\subparagraph}{0cm}{*1}{1em}

\setlength{\parindent}{0cm}
\setlength{\parskip}{0pt}

\usepackage{caption}
\captionsetup[table]{labelfont=bf,textfont=it,labelsep=newline,justification=raggedright,singlelinecheck=false}
\captionsetup[figure]{labelfont=bf,textfont=it,labelsep=newline,justification=raggedright,singlelinecheck=false}

\usepackage{xurl}
\tolerance=9999
\emergencystretch=3em
\hbadness=10000

\usepackage{bookmark}
\IfFileExists{xurl.sty}{\usepackage{xurl}}{} % add URL line breaks if available
\urlstyle{same}
\hypersetup{
  pdflang={es},
  hidelinks,
  pdfcreator={LaTeX via pandoc}}

\author{}
\date{}

\begin{document}

\setstretch{1}
\begin{center}

\textbf{REACTIVIDAD DE EXTREMO A EXTREMO}

\emph{D. A. Flores Cordero}

7690-14-745 Universidad Mariano Gálvez

047 Seminario de Tecnologías de Información

\emph{dfloresc3@miumg.edu.gt}

\end{center}

\section{Resumen}\label{resumen}

Este trabajo investiga las buenas prácticas para el desarrollo de
aplicaciones reactivas, con el objetivo de identificar los criterios que
determinan una implementación reactiva efectiva más allá de la simple
elección de un framework. La investigación partió de cuatro fuentes: una
comparación general entre programación reactiva e imperativa junto con
sus herramientas principales (RxJS, ReactiveX, Spring WebFlux), el texto
original del Manifiesto de Sistemas Reactivos con sus cuatro pilares
---responsividad, resiliencia, elasticidad y orientación a mensajes---,
un análisis complementario de ese mismo manifiesto, y una revisión
técnica de Spring WebFlux como implementación reactiva de backend en
Java. A partir de estas fuentes se identificaron tres prácticas
centrales: extender el enfoque no bloqueante a todas las capas de la
aplicación (controladores, acceso a datos, comunicación externa),
gestionar explícitamente el backpressure para no saturar al consumidor,
y aislar los fallos por componente para no comprometer el sistema
completo. Los resultados muestran que la reactividad parcial ---aplicar
el paradigma solo en una capa mientras el resto de la aplicación
permanece bloqueante--- anula buena parte de sus beneficios de
escalabilidad y eficiencia. Se concluye que la programación reactiva no
es una decisión de framework, sino una decisión arquitectónica de
extremo a extremo, y que su adopción debe evaluarse según la
concurrencia real que exige el sistema, ya que en aplicaciones pequeñas
puede introducir complejidad innecesaria sin un beneficio proporcional.

\textbf{Palabras claves:} programación reactiva, buenas prácticas,
backpressure, resiliencia, Spring WebFlux

\section{Desarrollo del tema}\label{desarrollo-del-tema}

La programación reactiva es un paradigma orientado a construir sistemas
que reaccionan de forma eficiente a flujos constantes de datos y eventos
mediante un enfoque no bloqueante (Quind, 2025). El Manifiesto de
Sistemas Reactivos formalizó esta idea al definir cuatro propiedades que
debe cumplir un sistema para considerarse reactivo: ser responsivo,
resiliente, elástico y orientado a mensajes (Bonér et al., 2014). Su
promesa ---mayor escalabilidad, mejor uso de recursos y menos problemas
de concurrencia--- solo se cumple cuando se aplican ciertas prácticas de
diseño; adoptar una librería reactiva sin cambiar la arquitectura
alrededor no basta.

\subsection{Más allá del framework: una decisión
arquitectónica}\label{muxe1s-alluxe1-del-framework-una-decisiuxf3n-arquitectuxf3nica}

Existe la tentación de tratar la reactividad como un detalle de
implementación: cambiar \texttt{RestTemplate} por \texttt{WebClient}, o
envolver una consulta en un \texttt{Observable}, y dar por hecho que la
aplicación ya es reactiva. Sin embargo, para obtener mejores beneficios
de un framework como Spring WebFlux es necesario que toda la aplicación
siga una arquitectura reactiva y no bloqueante en todas sus capas:
controladores, servicios, repositorios y comunicaciones externas (Quind,
2025). Una sola capa bloqueante en la cadena vuelve a introducir el
problema que la reactividad buscaba resolver.

\subsection{Extender el enfoque no bloqueante a todas las
capas}\label{extender-el-enfoque-no-bloqueante-a-todas-las-capas}

Esta primera práctica se traduce en decisiones concretas de
infraestructura. En el acceso a datos, corresponde usar bases de datos
compatibles con el paradigma ---MongoDB reactivo o bases relacionales
mediante R2DBC--- para evitar que una consulta bloquee un hilo mientras
espera respuesta (Quind, 2025). En la comunicación con servicios
externos, corresponde usar clientes HTTP reactivos como
\texttt{WebClient} en lugar de \texttt{RestTemplate}. La misma lógica
aplica en frontend: RxJS permite manejar los eventos de entrada del
usuario como flujos observables en vez de callbacks anidados,
integrándose con Angular, React o Vue (OBS Business School, s.f.).

\subsection{Gestionar el backpressure de forma
explícita}\label{gestionar-el-backpressure-de-forma-expluxedcita}

Procesar datos como flujo continuo trae un riesgo que la programación
imperativa no tiene: un productor más rápido que su consumidor. El
backpressure permite gestionar la velocidad con la que se envían y
procesan los datos para evitar saturaciones (Quind, 2025), por ejemplo
limitando el tamaño del buffer que un flujo puede acumular antes de
descartar o postergar nuevos elementos. Ignorar este mecanismo
---tratarlo como un detalle opcional del framework--- es una causa común
de que sistemas ``reactivos'' terminen consumiendo más memoria que su
equivalente bloqueante bajo carga alta.

\subsection{Aislar fallos: la resiliencia del Manifiesto
Reactivo}\label{aislar-fallos-la-resiliencia-del-manifiesto-reactivo}

El Manifiesto de Sistemas Reactivos define la resiliencia como la
capacidad de un sistema de permanecer responsivo frente a fallos,
lograda mediante replicación, contención, aislamiento y delegación: los
fallos se manejan dentro de cada componente, de forma que cualquier
parte del sistema pueda fallar y recuperarse sin comprometer el conjunto
(Bonér et al., 2014). Vilanou (s.f.) resume esta misma idea al señalar
que los fallos se trabajan de forma aislada para no afectar a los demás
componentes. En la práctica, esto significa diseñar cada componente para
fallar de forma contenida ---por ejemplo, un timeout en una llamada
externa no debería derribar el hilo principal de la aplicación--- en
lugar de asumir que un flujo de datos bien construido nunca falla.

\subsection{Elegir la herramienta según el lenguaje y el
paradigma}\label{elegir-la-herramienta-seguxfan-el-lenguaje-y-el-paradigma}

No existe una única herramienta reactiva: la elección depende del
entorno. En JavaScript, RxJS es la opción dominante para manejar eventos
y datos asíncronos mediante observables (OBS Business School, s.f.). En
Android, ReactiveX ---a través de RxJava y RxAndroid--- reduce el
acoplamiento entre componentes y optimiza peticiones asíncronas a APIs
(OBS Business School, s.f.). En backend Java, Spring WebFlux, construido
sobre Project Reactor con los tipos \texttt{Mono} y \texttt{Flux}, es la
alternativa reactiva a Spring MVC (Quind, 2025). Empresas como Netflix,
Alibaba y Tencent usan esta última para sostener millones de solicitudes
concurrentes (Quind, 2025).

\subsection{Cuándo no conviene la programación
reactiva}\label{cuuxe1ndo-no-conviene-la-programaciuxf3n-reactiva}

La última buena práctica es reconocer cuándo no aplicar el paradigma. En
aplicaciones pequeñas y simples, la programación reactiva puede ser
excesiva frente a la programación imperativa, que resulta más fácil de
implementar y mantener (OBS Business School, s.f.). Adoptarla sin una
necesidad real de concurrencia alta agrega complejidad de código y una
curva de aprendizaje ---flujos, operadores, backpressure--- sin un
beneficio proporcional.

\section{Observaciones y comentarios}\label{observaciones-y-comentarios}

Investigar tres fuentes con enfoques distintos ---una comparativa
general, el Manifiesto Reactivo y un caso técnico de Spring WebFlux---
dejó claro que la mayoría de la documentación disponible describe
\emph{qué} es la programación reactiva, pero muy poca detalla
\emph{cómo} aplicarla de forma consistente en todas las capas. El mayor
hallazgo no fue una herramienta específica, sino confirmar que la
resiliencia y el backpressure suelen tratarse como temas secundarios
frente al atractivo de la sintaxis de flujos, cuando en la práctica son
los que sostienen el sistema bajo carga real.

\section{Conclusiones}\label{conclusiones}

\begin{enumerate}
\def\labelenumi{\arabic{enumi}.}
\tightlist
\item
  La programación reactiva es una decisión arquitectónica de extremo a
  extremo, no la elección de un framework o una librería puntual.
\item
  Aplicar el paradigma en una sola capa mientras el resto de la
  aplicación permanece bloqueante anula buena parte de sus beneficios de
  escalabilidad.
\item
  Gestionar el backpressure y aislar los fallos por componente son
  prácticas tan necesarias como el manejo de flujos asíncronos, aunque
  reciben menos atención en la documentación general.
\item
  La herramienta correcta depende del entorno: RxJS en frontend
  JavaScript, ReactiveX en Android y Spring WebFlux en backend Java.
\item
  La programación reactiva no conviene en aplicaciones pequeñas sin
  necesidad real de alta concurrencia, donde su complejidad adicional no
  se justifica.
\end{enumerate}

\section{Bibliografía}\label{bibliografuxeda}

Bonér, J., Farley, D., Kuhn, R., \& Thompson, M. (2014). \emph{El
Manifiesto de Sistemas Reactivos} (v2.0).
https://www.reactivemanifesto.org/es

OBS Business School. (s.f.). \emph{Cómo aplicar la programación reactiva
en aplicaciones web modernas}.
https://www.obsbusiness.school/blog/como-aplicar-la-programacion-reactiva-en-aplicaciones-web-modernas-cp

Quind. (2025). \emph{Programación reactiva con Spring WebFlux:
eficiencia y escalabilidad para el backend}.
https://quind.io/blog/desarrollo/programacion-reactiva-con-spring-webflux-eficiencia-y-escalabilidad-para-el-backend/

Vilanou, R. (s.f.). \emph{¿Qué es la programación reactiva?}. Doonamis.
https://www.doonamis.com/que-es-la-programacion-reactiva/

\begin{center}\rule{0.5\linewidth}{0.5pt}\end{center}

Repositorio del documento (código fuente LaTeX):
\url{https://github.com/dan-desa/seminario-tecnologias-2026/tree/main/foro-02}

\end{document}
